% =====================================================================
%  21.4  Persistent roots planted on a tree: k-th / count on a u--v path
%  Exports: sec:c19-ontrees  def:c19-rootpath  prop:c19-fourroots
%           lst:c19-persist-tree  lst:c19-pathkth  lst:c19-cot-soln
%           prob:c19-cot  fig:c19-path
%  Runs on the SHARED 10-node tree T of Ch.12-17 (master lines 74-76):
%    edges 1-2,1-3,2-4,2-5,3-6,3-7,5-8,7-9,9-10 ; root of T = node 1.
%    weights w(node 1..10) = (3,2,4,5,1,1,5,3,7,0).
%    four-root combine root_u + root_v - root_lca - root_parent(lca):
%      path 3--8 = 3-1-2-5-8, lca=1, parent(lca)=none.
%        root_3={3,4} root_8={3,2,1,3} root_1={3} => multiset {1,2,3,3,4}
%        count<=3 = 4.   (matches 00_opener line 34-35 and master line 47-48)
%      path 8--10 = 8-5-2-1-3-7-9-10, lca=1.
%        root_8={3,2,1,3} root_10={3,4,5,7,0} root_1={3} => {0,1,2,3,3,4,5,7}
%        k-th (k=6) = 4. (matches 00_opener line 43 and master line 48)
%  SPOJ COT (live-verified on spoj.com/problems/COT): "Count on a tree".
%    N,M<=100000; each node an integer weight; query u v k = k-th MINIMUM
%    weight on the path u--v.  Official sample:
%      8 5 / w=(105,2,9,3,8,5,7,7) / edges 1-2,1-3,1-4,3-5,3-6,3-7,4-8
%      (2,5,1)->2  (2,5,2)->8  (2,5,3)->9  (2,5,4)->105  (7,8,2)->7.
%      Hand-replay: path 2-5=2-1-3-5 weights{2,105,9,8}sorted{2,8,9,105};
%      path 7-8=7-3-1-4-8 weights{7,9,105,3,7}sorted{3,7,7,9,105} 2nd=7. OK.
% =====================================================================
\section{Persistent Roots on a Tree: \(k\)-th and Count on a Path}\label{sec:c19-ontrees}

\begin{previouslybox}[Ch.~20's lastword promised ``reads against past versions''; here the versions live \emph{on the tree}]
\(\secref{c19-chairman}\) made the versions the array's prefixes so that a \emph{subarray} was a two-root difference. Chapter~20, closing on persistent structures, promised this chapter would answer ``\(k\)-th inside a static range \emph{and reads against past versions}.'' The tree is the next static range: give each node \(x\) a value, root the tree at some \(s\), and ask about the multiset of values on the path \(u\text{--}v\). The trick is to keep planting one persistent root per \term{tree node} rather than per array index: \(\mathit{root}_x\) is the frequency tree, over the value axis, of the values on the root-to-\(x\) path. Then a path is not a difference of \emph{two} roots but of \emph{four}, because two root-to-node chains overlap along their common ancestry up to the lowest common ancestor. This is the same path-copying spine of \(\secref{c19-copy}\) and the same two-root \(k\)-th of \(\secref{c19-chairman}\), with Ch.~11's LCA bolted on to name the two extra roots.
\end{previouslybox}

\begin{topicintro}
Root the tree \(T\) at node \(s=1\). DFS from the root and, on entering a node \(v\) with parent \(p\), set \(\mathit{root}_v=\texttt{insert}(\mathit{root}_p,\rho_v)\) -- one path copy of the parent's frequency root, adding node \(v\)'s compressed value \(\rho_v\). So \(\mathit{root}_v\) holds the whole \term{root-to-\(v\) path multiset} and, exactly as in \(\secref{c19-chairman}\), shares every node \(\mathit{root}_p\) had that this insert did not touch.

Now the path \(u\text{--}v\). Write \(\ell=\mathrm{lca}(u,v)\). The two chains root\(\to u\) and root\(\to v\) meet at \(\ell\) and coincide above it. Adding \(\mathit{root}_u+\mathit{root}_v\) counts the nodes on root\(\to\ell\) \emph{twice} and \(\ell\) itself \emph{twice}; the path \(u\text{--}v\) should contain \(\ell\) once and none of its strict ancestors. Subtracting \(\mathit{root}_\ell\) removes one copy of the shared spine \emph{including} \(\ell\) -- which overshoots by deleting \(\ell\) entirely -- so we add \(\ell\) back by subtracting only \(\mathit{root}_{\mathrm{parent}(\ell)}\) instead. The clean statement:
\[
  \text{path } u\text{--}v \;=\; \mathit{root}_u \;+\; \mathit{root}_v \;-\; \mathit{root}_{\ell} \;-\; \mathit{root}_{\mathrm{parent}(\ell)},
  \qquad \ell=\mathrm{lca}(u,v),
\]
read as a per-value-count combination of \emph{four} frequency roots. Every range statistic \(\secref{c19-chairman}\) could take from two roots -- a count \(\le x\), a \(k\)-th by joint descent -- takes the analogous four-root combination here, descending four roots in lockstep.

\textbf{Prerequisites.} \(\secref{c19-copy}\) (\(\defref{c19-pathcopy}\), the node pool) and \(\secref{c19-chairman}\) (\(\coderef{c19-chairman}\)'s \texttt{insert}, \(\coderef{c19-kth}\)'s joint descent) are assumed and \emph{reused}, not re-derived. Ch.~11 (\chref{c11}) owns binary lifting and \(\mathrm{lca}()\); Ch.~5 (\chref{c24}) owns the \(tin/tout\) DFS order (used only to say what a subtree flattening is). \(\secref{c19-sparse}\) closes the chapter's engine map.

\textbf{After this section you can}
\begin{itemize}[nosep]
  \item plant a persistent frequency root on each tree node with one \code{insert} per DFS step; \hfill\lv{3}--\lv{4}
  \item combine four roots (\(u,v,\ell,\mathrm{parent}(\ell)\)) for any path multiset query; \hfill\lv{4}
  \item answer the \(k\)-th smallest value on a path \(u\text{--}v\) by a \(\bigO{\log V}\) four-root descent; \hfill\lv{4}
  \item explain why the `extra' root is \(\mathrm{parent}(\ell)\) and not \(\ell\) (the LCA-inclusion off-by-one). \hfill\lv{4}
\end{itemize}
\end{topicintro}

\begin{recognitionbox}[Pattern card -- ``\(k\)-th / count of node-values along a tree path, static tree'']
\reckey{a \emph{static} weighted tree; queries ``\(k\)-th smallest value on the path \(u\text{--}v\)'' / ``how many nodes on \(u\text{--}v\) have value \(\le x\)'' / ``count colour \(c\) on a path''; node values on a tree, many path order-statistic questions, no updates -- SPOJ~COT is the archetype.}
\recuse{
  \item \(k\)-th node-value on a path, static tree, \(N\le10^5\) -- this section's four-root descent (\(\coderef{c19-pathkth}\)), SPOJ~COT;
  \item count of a value / a threshold on a path -- the same four roots, summed;
  \item a `path aggregate that is a frequency' -- plant a root per node and combine at the LCA.
}
\recnotuse{
  \item the values sit on a flat \emph{subarray}, not a tree path -- that is \(\secref{c19-chairman}\) with two roots;
  \item the tree is \emph{updated} (repaint paths, add nodes) -- persistence goes stale; that is HLD + lazy, Ch.~22 (\chref{c16});
  \item the query is a \emph{subtree} rather than a path -- flatten with \(tin/tout\) (Ch.~5, \chref{c24}) and the two-root prefix trick already does it.
}
One-line test: \emph{a static tree, node values, and ``which value sits \(k\)-th along the path \(u\text{--}v\)'' -- plant a frequency root per node and subtract up from four of them.}
\end{recognitionbox}

% ---------------------------------------------------------------
\subsection{Intuition: two chains share a spine, so four roots cancel it}
% ---------------------------------------------------------------

A path \(u\text{--}v\) is the concatenation of \(u\)-up-to-\(\ell\) and \(\ell\)-down-to-\(v\), meeting at \(\ell=\mathrm{lca}(u,v)\). The persistent roots give us, for any node, the \emph{whole} root-to-node chain in one tree. But \(u\) and \(v\) both descend from \(s\), so \(\mathit{root}_u\) and \(\mathit{root}_v\) each already contain the spine \(s\to\ell\). If we simply add them, every node on \(s\to\ell\) is counted twice, and \(\ell\) is counted twice, while the path should keep only \(\ell\) once and none of its ancestors.

Subtracting \(\mathit{root}_\ell\) removes exactly one copy of \(s\to\ell\) \emph{including} \(\ell\). That leaves \(\ell\) counted once from the surviving \(\mathit{root}_v\)/\(\mathit{root}_u\) contribution minus the deletion... which is one too few: \(\ell\) ends up removed. The fix is to delete the spine only \emph{above} \(\ell\), i.e.\ subtract \(\mathit{root}_{\mathrm{parent}(\ell)}\) -- \(\mathit{root}_\ell\) without \(\ell\)'s own value. Hence the four roots \(u,v,\ell,\mathrm{parent}(\ell)\). It is the same inclusion--exclusion you would write for `size of a path set,' dressed in frequency trees.

Once the combination is set, everything transfers verbatim from \(\secref{c19-chairman}\): a count \(\le x\) on the path is
\[
  \bigl(T[\mathit{root}_u].\mathit{cnt}+T[\mathit{root}_v].\mathit{cnt}-T[\mathit{root}_\ell].\mathit{cnt}-T[\mathit{root}_{\mathrm{parent}(\ell)}].\mathit{cnt}\bigr)
\]
accumulated over the value prefix \([1..t]\); and the \(k\)-th is a \emph{four-root joint descent} -- at each value level compute the left-child difference of all four and let it decide the branch. \(\mathrm{lca}\) and \(\mathrm{parent}\) come from Ch.~11's binary lifting; building the root chain is one \code{insert} per edge, folded into the same DFS that computes \(\mathit{depth}\) and the lifting table.

\begin{definition}{A persistent root per tree node}{c19-rootpath}
Root \(T\) at \(s\). Let \(\rho_v\) be the compressed value rank at node \(v\). Assign \(\mathit{root}_v\) for every node by DFS: \(\mathit{root}_s=\texttt{insert}(\mathbf{0},\rho_s)\) (from the null tree) and, for a child \(c\) of \(v\), \(\mathit{root}_c=\texttt{insert}(\mathit{root}_v,\rho_c)\), where \texttt{insert} is the frequency path-copy of \(\coderef{c19-chairman}\). Thus \(\mathit{root}_v\)'s tree holds, for each value interval, the number of nodes on the root-to-\(v\) path whose value rank lies there; the total node cost is \(\bigO{N\log V}\).
\end{definition}

\begin{proposition}{The four-root path combination}{c19-fourroots}
For a static rooted tree with a persistent frequency root per node (\(\defref{c19-rootpath}\)) and \(\ell=\mathrm{lca}(u,v)\), the multiset of value ranks on the simple path \(u\text{--}v\) (endpoints included) is, at every value node,
\[
  \mathit{cnt}(u,v\ \text{path}) \;=\; \mathit{cnt}(\mathit{root}_u)+\mathit{cnt}(\mathit{root}_v)-\mathit{cnt}(\mathit{root}_{\ell})-\mathit{cnt}(\mathit{root}_{\mathrm{parent}(\ell)}).
\]
Consequently a count \(\le x\) on the path and the path's \(k\)-th smallest are each answered in \(\bigO{\log V}\), the latter by a four-root joint descent, and every root-to-node value is stored once so the whole structure is \(\bigO{N\log V}\) nodes with \emph{no} updates.
\end{proposition}
\begin{proofbox}
Every node \(w\) on the path \(u\text{--}v\) lies on exactly one of the two chains \(u\to\ell\) (inclusive of \(u,\ell\)) or \(\ell\to v\) excluding \(\ell\). The root-to-\(u\) chain = \(s\to\ell\) plus \(\ell\to u\); the root-to-\(v\) chain = \(s\to\ell\) plus \(\ell\to v\). Their sum counts each node of \(s\to\ell\) twice and each node strictly below \(\ell\) on either side once. The path wants the below-\(\ell\) nodes (kept, once) and \(\ell\) (kept once), but none of \(s\to\)parent\((\ell)\) (dropped entirely). \(\mathit{root}_\ell\) counts \(s\to\ell\) once; \(\mathit{root}_{\mathrm{parent}(\ell)}\) counts \(s\to\)parent\((\ell)\) once; their difference is \(\{\ell\}\). Subtracting \(\mathit{root}_\ell\) and adding back... i.e.\ subtracting the pair removes one copy of \(s\to\)parent\((\ell)\) \emph{and} leaves \(\ell\)'s single copy intact: nodes of \(s\to\)parent\((\ell)\) had count \(2-1-1=0\); \(\ell\) has \(2-1-0=1\); below-\(\ell\) path nodes have \(1-0-0=1\); off-path nodes have \(0\). Each per-value count is therefore exact, and a joint descent that compares the four-way left-child difference to \(k\) is the identical argument to \thmref{c19-kth}, so the \(k\)-th leaf is reached in \(\bigO{\log V}\). The build visits each node once and does one \(O(\log V)\) insert per node.
\end{proofbox}

\noindent\textbf{In plain words.}  Photograph the shelf (the value histogram) at \emph{every} node, showing everything from the root down to that node -- each photo is one cheap edit of its parent's. To see a path \(u\text{--}v\), add \(u\)'s and \(v\)'s photos; the shared top part (root down to the meeting point \(\ell\)) appears twice, so subtract the meeting point's photo and its parent's photo to cancel that top and leave \(\ell\) exactly once. Four photos, and what is left is the path. Everything the two-root version could do (count, \(k\)-th) works here with four.

% ---------------------------------------------------------------
\subsection{Planting the roots, and the four-root \(k\)-th}
% ---------------------------------------------------------------

\(\coderef{c19-persist-tree}\) folds the chairman \texttt{insert} into the tree DFS; \(\coderef{c19-pathkth}\) turns \(\coderef{c19-kth}\)'s two-way descent four-way.

\begin{lstlisting}[caption={plant one frequency root per tree node: each DFS step path-copies the parent root and inserts the node's own value rank}, label={lst:c19-persist-tree}]
int ROOTV[MAXN];               // ROOTV[v] = frequency tree over values on path (root .. v)
void dfs(int v, int p){
    ROOTV[v] = insert(ROOTV[p], 1, V, rk[v]);   // one path copy of parent's root, add node v
    for (int to : adj[v]) if (to != p) dfs(to, v);
}
// call: ROOTV[0] = 0 (the null tree);  dfs(s, 0)  with s the chosen tree root.
// insert / node pool are exactly those of Code lst:c19-chairman -- nothing new.
\end{lstlisting}

\begin{lstlisting}[caption={k-th smallest value on the path u--v: descend four roots at once, the LCA and its parent being the two subtracted ones}, label={lst:c19-pathkth}]
// a=ROOTV[u], b=ROOTV[v], c=ROOTV[lca], d=ROOTV[parent(lca)]; value interval [tl..tr]
int kthOnPath(int a, int b, int c, int d, int tl, int tr, int k){
    if (tl == tr) return tl;                     // leaf rank = answer's compressed value
    int tm = (tl + tr) / 2;
    // # path values landing in the LOWER half [tl..tm], from four roots:
    int leftCount = T[T[a].l].cnt + T[T[b].l].cnt
                  - T[T[c].l].cnt - T[T[d].l].cnt;
    if (k <= leftCount)
        return kthOnPath(T[a].l, T[b].l, T[c].l, T[d].l, tl, tm, k);
    else
        return kthOnPath(T[a].r, T[b].r, T[c].r, T[d].r, tm+1, tr, k - leftCount);
}
// lca(u,v) and parent(lca) come from Ch.11 binary lifting; pass ROOTV of all four.
\end{lstlisting}

A count on a path is the same four-way subtraction but summed to the value prefix \([1..t]\): descend to node \([1..t]\) once and read \(T[a].\mathit{cnt}+T[b].\mathit{cnt}-T[c].\mathit{cnt}-T[d].\mathit{cnt}\). The only new arithmetic versus \(\secref{c19-chairman}\) is two more terms -- and two more pointers to carry down the descent.

% ---------------------------------------------------------------
\subsection{The hand trace on the shared tree \(T\)}
% ---------------------------------------------------------------

\begin{figure}[H]
\centering
\begin{tikzpicture}[diagram, scale=0.6, transform shape]
  \node[boxnode, primaryblue, minimum size=8mm, font=\scriptsize] (n1) at (3,3) {\(1\)};
  \node[boxnode, primaryblue, minimum size=7mm, font=\scriptsize] (n2) at (1.5,1.9) {\(2\)};
  \node[boxnode, primaryblue, minimum size=7mm, font=\scriptsize] (n3) at (4.5,1.9) {\(3\)};
  \node[boxnode, gray, minimum size=6mm, font=\scriptsize] (n4) at (0.6,0.8) {\(4\)};
  \node[boxnode, primaryblue, minimum size=6mm, font=\scriptsize] (n5) at (2.0,0.8) {\(5\)};
  \node[boxnode, gray, minimum size=6mm, font=\scriptsize] (n6) at (3.9,0.8) {\(6\)};
  \node[boxnode, primaryblue, minimum size=6mm, font=\scriptsize] (n7) at (5.3,0.8) {\(7\)};
  \node[boxnode, difforange, minimum size=6mm, font=\scriptsize] (n8) at (2.0,-0.3) {\(8\)};
  \node[boxnode, primaryblue, minimum size=6mm, font=\scriptsize] (n9) at (5.3,-0.3) {\(9\)};
  \node[boxnode, difforange, minimum size=6mm, font=\scriptsize] (n10) at (5.3,-1.3) {\(10\)};
  \draw[flow] (n1)--(n2); \draw[flow] (n1)--(n3);
  \draw[flow] (n2)--(n4); \draw[flow] (n2)--(n5); \draw[flow] (n3)--(n6); \draw[flow] (n3)--(n7);
  \draw[flow] (n5)--(n8); \draw[flow] (n7)--(n9); \draw[flow] (n9)--(n10);
  \node[font=\tiny\ttfamily, text=darkgray, anchor=west] at (3.3,3) {\(w{=}3\)};
  \node[font=\tiny\ttfamily, text=darkgray, anchor=west] at (1.75,1.9) {\(2\)};
  \node[font=\tiny\ttfamily, text=darkgray, anchor=west] at (4.75,1.9) {\(4\)};
  \node[font=\tiny\ttfamily, text=darkgray, anchor=west] at (2.25,0.8) {\(1\)};
  \node[font=\tiny\ttfamily, text=darkgray, anchor=west] at (5.55,0.8) {\(5\)};
  \node[font=\tiny\ttfamily, text=orange!60!black, anchor=west] at (2.25,-0.3) {\(3\)};
  \node[font=\tiny\ttfamily, text=orange!60!black, anchor=west] at (5.55,-1.3) {\(0\)};
  \node[font=\scriptsize, text=orange!55!black, anchor=north] at (3,-2.1) {path 8--10: 8-5-2-1-3-7-9-10,  lca(8,10)=1};
\end{tikzpicture}
\caption{The shared 10-node tree \(T\) (root \(1\)) with node weights \(w=(3,2,4,5,1,1,5,3,7,0)\).  The highlighted spine is the query path \(8\text{--}10\).  The path multiset is \(\mathit{root}_8+\mathit{root}_{10}-\mathit{root}_1-\mathit{root}_{\mathrm{parent}(1)}\); since \(\mathrm{lca}(8,10)=1\) and node \(1\) has no parent, the fourth root is the null tree.}
\label{fig:c19-path}
\end{figure}

\begin{tracebox}[k-th on path 8 to 10 with k equals 6 gives 4, combining four roots]
\textbf{Roots (root-to-node chains).} \(\mathit{root}_8\): nodes \(1,2,5,8\) values \(\{3,2,1,3\}\). \(\mathit{root}_{10}\): nodes \(1,3,7,9,10\) values \(\{3,4,5,7,0\}\). \(\mathit{root}_1\) (lca): \(\{3\}\). \(\mathit{root}_{\mathrm{parent}(1)}=\mathbf{0}\).\\
\textbf{Combine.} \(\{3,2,1,3\}+\{3,4,5,7,0\}-\{3\}-\varnothing=\{0,1,2,3,3,4,5,7\}\) -- exactly the eight nodes \(8,5,2,1,3,7,9,10\).\\
\textbf{Fork 1, value interval \([0..7]\), lower half \([0..3]\):} count \(=5\) (the values \(0,1,2,3,3\)).  \(k=6>5\) \(\Rightarrow\) go upper \([4..7]\) with \(k'=6-5=1\).\\
\textbf{Fork 2, \([4..7]\), lower \([4..5]\):} count \(=2\) (values \(4,5\)).  \(k'=1\le2\) \(\Rightarrow\) stay in \([4..5]\).\\
\textbf{Fork 3, \([4..5]\), lower \([4]\):} count \(=1\).  \(k'=1\le1\) \(\Rightarrow\) value \(\boxed{4}\).\\
\textbf{Cross-check by sorting:} path \(=\{0,1,2,3,3,4,5,7\}\); the 6th smallest is \(4\).  Three four-way reads, and the meeting node \(1\) (weight \(3\)) is counted once, as \propref{c19-fourroots} promises.
\end{tracebox}

\begin{tracebox}[count values <= 3 on path 3 to 8 -- the four-root sum is 4]
\textbf{Roots.} \(\mathit{root}_3\): \(1,3\) values \(\{3,4\}\).  \(\mathit{root}_8\): \(1,2,5,8\) values \(\{3,2,1,3\}\).  \(\mathit{root}_{\mathrm{lca}(3,8)}=\mathit{root}_1=\{3\}\); \(\mathit{root}_{\mathrm{parent}(1)}=\varnothing\).\\
\textbf{Combine:} \(\{3,4\}+\{3,2,1,3\}-\{3\}=\{1,2,3,3,4\}\) -- the five nodes \(3,1,2,5,8\).\\
\textbf{Count \(\le 3\):} \(1,2,3,3\) \(\Rightarrow\) \(\boxed{4}\).\\
\textbf{Cross-check} by brute walk of the path \(3\text{--}8=3\text{-}1\text{-}2\text{-}5\text{-}8\) with weights \(4,3,2,1,3\): entries \(\le3\) are \(3,2,1,3\) -- four.  The \(00\_opener\) figure path \(3\text{--}8=(1,2,3,3,4)\) and answer \(4\) match.
\end{tracebox}

\begin{complexitybox}[Persistent roots on a tree]
\textbf{Build} \(\bigO{N\log V}\) nodes and time (one \code{insert} per node inside the DFS that also fills LCA lifting). \textbf{Query} \(\bigO{\log V}\) for a path \(k\)-th or count (four-root descent), plus \(\bigO{\log N}\) to get \(\mathrm{lca}\) (Ch.~11). \textbf{Memory} \(\bigO{N\log V}\) nodes: with \(N\le10^5,V\le10^5\) budget \(\sim1.7\cdot10^6\) nodes -- size the pool as \(\thmref{c19-memory}\) warns. \textbf{Static}: a value repaint invalidates every later root on that subtree -- move to HLD/lazy (Ch.~22 \chref{c16}).
\end{complexitybox}

% ---------------------------------------------------------------
\subsection{Problem Lab: SPOJ COT -- the \(k\)-th weight on a tree path}
% ---------------------------------------------------------------

\solvedproblem{c19-cot}{Count on a tree}{SPOJ COT}
  {https://www.spoj.com/problems/COT/}{SPOJ L6}
  {chairman roots planted per tree node; a path is four roots; \(k\)-th by joint four-way descent}

\begin{problemstatement}[SPOJ COT Count on a tree]
A tree with \(N\) nodes numbered \(1..N\), each node carrying an integer weight. Answer \(M\) queries of the form \(u\ v\ k\): the \(k\)-th \emph{minimum} weight on the path from node \(u\) to node \(v\) (endpoints included).

\textbf{Input.} Line 1: \(N\) and \(M\) (\(N,M\le 100000\)). Line 2: the \(N\) weights. Next \(N-1\) lines: edges \(u\ v\). Next \(M\) lines: queries \(u\ v\ k\).

\textbf{Output.} For each query, its answer, one per line.

\textbf{Example (official).}
\begin{verbatim}
Input:
8 5
105 2 9 3 8 5 7 7
1 2
1 3
1 4
3 5
3 6
3 7
4 8
2 5 1
2 5 2
2 5 3
2 5 4
7 8 2

Output:
2
8
9
105
7
\end{verbatim}
\end{problemstatement}

\begin{attemptgate}[Attempt it yourself first -- \lv{4}, give yourself 35 minutes]
A path is not a prefix, so \(\secref{c19-chairman}\)'s two roots are not enough on their own -- but root-to-node chains are exactly `tree prefixes.' Plant a chairman root per node, then reach a path through the LCA. The subtlety to find: which nodes do you subtract so \(\ell\) is counted \emph{once}?
\peek{why can \(\mathit{root}_u+\mathit{root}_v\) not stand alone for the path \(u\text{--}v\)?}
\nudge{the two chains share the spine root\(\to\ell\); it is counted twice and must go. Removing \(\mathit{root}_\ell\) also deletes \(\ell\), so remove \(\mathit{root}_{\mathrm{parent}(\ell)}\) instead -- that keeps \(\ell\) once.}
\fullhint{Compress the weights to ranks; DFS from node~1 building \(\mathit{root}_v=\texttt{insert}(\mathit{root}_{\mathrm{parent}(v)},\rho_v)\) and the binary-lifting table (Ch.~11). For query \((u,v,k)\): \(\ell=\mathrm{lca}(u,v)\), \(d=\mathrm{parent}(\ell)\); run the four-root descent \(t=\mathit{kthOnPath}(\mathit{root}_u,\mathit{root}_v,\mathit{root}_\ell,\mathit{root}_d,1,V,k)\); print \(\mathit{vals}[t]\). Total \(\bigO{(N+M)\log N}\approx(2\cdot10^5)\cdot17\approx3.4\cdot10^6\) node steps, \ac. Watch the null root when \(\ell\) is the tree root.}
\end{attemptgate}

\solstep{Understanding}
\(N,M\le10^5\), static weighted tree, path \(k\)-th. Persistent frequency roots turn each root-to-node chain into one tree; a path is a four-root combination (\(\propref{c19-fourroots}\)); the \(k\)-th is a joint four-way descent. Coordinate compression handles weights like \(105\) (and duplicates like the two \(7\)s, which share a leaf).

\solstep{Key Observations}
\begin{observation}
The `path multiset' is inclusion--exclusion over \emph{root-to-node chains}: \(u\) and \(v\) chains, minus the shared spine counted twice. That forces the LCA \emph{and its parent} into the formula -- the two roots \(\secref{c19-chairman}\) never needed.
\end{observation}
\begin{observation}
\(\ell\)'s own value is on the path exactly once. Subtracting \(\mathit{root}_\ell\) deletes it, so you subtract \(\mathit{root}_{\mathrm{parent}(\ell)}\) to spare it -- this single choice is the whole off-by-one risk of the section.
\end{observation}
\begin{observation}
The root \(s\) has no parent: treat \(\mathit{root}_{\mathrm{parent}(s)}=\mathit{root}_0\) as the null (all-zero) tree so the formula is uniform even when \(\ell=s\).
\end{observation}

\solstep{Think It Through With Me}
\begin{thinkbox}
\deadend{Two roots only, \(\mathit{root}_u-\mathit{root}_{l-1}\) style}
{There is no `index \(l-1\)' on a tree; \(\mathit{root}_u\) and \(\mathit{root}_v\) overlap on the shared spine, and a bare difference does not isolate the path. You need the LCA to name where they diverge.}
\deadend{Subtract just \(\mathit{root}_\ell\)}
{\(\mathit{root}_u+\mathit{root}_v-\mathit{root}_\ell\) deletes \(\ell\) entirely (\(\ell\) appeared twice in the sum, once in the subtraction), so the path loses its meeting node -- off by one value, wrong for every query whose answer is \(\ell\)'s weight.}
\aha{Subtract \(\mathit{root}_\ell\) \emph{and} \(\mathit{root}_{\mathrm{parent}(\ell)}\): the difference of those two is exactly \(\{\ell\}\), so the spine vanishes while \(\ell\) stays once. Four roots, one clean descent.}
\end{thinkbox}

\solstep{The Approach Ladder}
\begin{center}\renewcommand{\arraystretch}{1.25}\setlength{\tabcolsep}{3pt}
\begin{tabular}{@{}c p{0.22\linewidth} p{0.26\linewidth} l c@{}}
\toprule
\textbf{\#} & \textbf{Approach} & \textbf{Idea} & \textbf{Time} & \textbf{Verdict}\\
\midrule
1 & \lv{2}\ walk + sort & collect path values, sort, index \(k\) & \(\bigO{M\,N\log N}\) & \tle \\
2 & \lv{3}\ HLD + merge-sort & decompose path, count per chain & \(\bigO{M\log^3 N}\) & \ac \\
3 & \lv{4}\ \textbf{chairman on tree} & four-root path, joint descent (\(\coderef{c19-pathkth}\)) & \(\bigO{(N+M)\log N}\) & \ac \\
4 & \lv{4}\ Mo on tree & offline path queries & \(\bigO{N\sqrt N}\) & \ac\\
\bottomrule
\end{tabular}\end{center}
Rung~2 (heavy-light, Ch.~22 \chref{c16}) works but re-pays a log; rung~4 (Mo's, Ch.~27 \chref{c21}) is offline and cannot answer in order. Rung~3 answers online in one \(\bigO{\log N}\) descent -- the point of planting the roots.

\solstep{Solution}
Compress weights to ranks \(1..V\). DFS from node \(1\): build each node's frequency root by one \code{insert} of its own rank on its parent's root, and fill the LCA lifting table. Per query, \(\ell=\mathrm{lca}(u,v)\), \(d=\) parent\((\ell)\) or \(0\); run the four-root \(k\)-th descent and un-compress the leaf.

\solstep{Proof of Correctness}
\begin{proofbox}
\textbf{Roots exact.} By \(\defref{c19-rootpath}\) and \(\thmref{c19-memory}\), \(\mathit{root}_v\) holds the root-to-\(v\) multiset after one path-copy insert per node, sharing all untouched nodes with \(\mathit{root}_{\mathrm{parent}(v)}\).\\
\textbf{Combination exact.} \(\propref{c19-fourroots}\)' per-node count check gives \(0\) for spine-above-\(\ell\), \(1\) for \(\ell\), \(1\) for below-\(\ell\) path nodes, \(0\) off-path -- i.e.\ exactly the path multiset. The four-root descent compares this difference to \(k\) at each value level; the same invariant as \thmref{c19-kth} lands on the \(k\)-th rank leaf. Printing \(\mathit{vals}[t]\) is the inverse compression.\\
\textbf{Complexity.} Build \(\bigO{N\log N}\); query \(\bigO{\log N}\) plus \(\bigO{\log N}\) for LCA. With \(N,M\le10^5\) total \(\approx3.4\cdot10^6\) steps -- \ac; node pool \(\sim N(\lceil\log_2 N\rceil+1)\approx1.8\cdot10^6\).
\end{proofbox}

\solstep{Implementation}
\begin{lstlisting}[caption={SPOJ COT full solution: per-node chairman roots + Ch.11 LCA, answer each query by a four-root joint descent}, label={lst:c19-cot-soln}]
#include <bits/stdc++.h>
using namespace std;
const int MAXN = 100005, LOGK = 17;
struct Node { int cnt, l, r; };
vector<Node> T;
int newnode(int c,int l,int r){ T.push_back({c,l,r}); return (int)T.size()-1; }
vector<int> adj[MAXN];
int wt[MAXN], rk[MAXN], ROOTV[MAXN];
int up[LOGK][MAXN], dep[MAXN];
vector<int> vals; int V;

int insert(int prev, int tl, int tr, int pos){
    int cur = newnode(T[prev].cnt+1, T[prev].l, T[prev].r);
    if (tl != tr){
        int tm=(tl+tr)/2;
        if (pos<=tm){ int nl=insert(T[prev].l,tl,tm,pos); T[cur].l=nl; }
        else        { int nr=insert(T[prev].r,tm+1,tr,pos); T[cur].r=nr; }
    }
    return cur;
}
int kthOnPath(int a,int b,int c,int d,int tl,int tr,int k){
    if (tl==tr) return tl;
    int tm=(tl+tr)/2;
    int lc = T[T[a].l].cnt + T[T[b].l].cnt - T[T[c].l].cnt - T[T[d].l].cnt;
    if (k<=lc) return kthOnPath(T[a].l,T[b].l,T[c].l,T[d].l,tl,tm,k);
    else       return kthOnPath(T[a].r,T[b].r,T[c].r,T[d].r,tm+1,tr,k-lc);
}
void dfs(int v,int p){
    up[0][v]=p; dep[v]=p?dep[p]+1:1;
    ROOTV[v]=insert(ROOTV[p],1,V,rk[v]);
    for (int j=1;j<LOGK;j++) up[j][v]=up[j-1][ up[j-1][v] ];
    for (int to: adj[v]) if (to!=p) dfs(to,v);
}
int lca(int a,int b){
    if (dep[a]<dep[b]) swap(a,b);
    int diff=dep[a]-dep[b];
    for (int j=0;j<LOGK;j++) if (diff&(1<<j)) a=up[j][a];
    if (a==b) return a;
    for (int j=LOGK-1;j>=0;j--) if (up[j][a]!=up[j][b]){ a=up[j][a]; b=up[j][b]; }
    return up[0][a];
}
int main(){
    int N,M; scanf("%d%d",&N,&M);
    vals.reserve(N);
    for (int i=1;i<=N;i++){ scanf("%d",&wt[i]); vals.push_back(wt[i]); }
    sort(vals.begin(),vals.end());
    vals.erase(unique(vals.begin(),vals.end()),vals.end()); V=vals.size();
    for (int i=1;i<=N;i++) rk[i]=lower_bound(vals.begin(),vals.end(),wt[i])-vals.begin()+1;
    for (int i=1;i<N;i++){ int u,v; scanf("%d%d",&u,&v); adj[u].push_back(v); adj[v].push_back(u); }
    T={ {0,0,0} }; ROOTV[0]=0; dep[0]=0; up[0][0]=0;
    dfs(1,0);
    while (M--){
        int u,v,k; scanf("%d%d%d",&u,&v,&k);
        int l=lca(u,v), d=up[0][l];
        int t=kthOnPath(ROOTV[u],ROOTV[v],ROOTV[l],ROOTV[d],1,V,k);
        printf("%d\n", vals[t-1]);
    }
}
\end{lstlisting}

\solstep{Trace on the Official Sample}
\begin{tracebox}[COT: path 2 to 5 and path 7 to 8]
Weights \(w=(105,2,9,3,8,5,7,7)\); tree \(1\{2,3,4\}\), \(3\{5,6,7\}\), \(4\{8\}\).\\
\textbf{Query (2,5,1..4).} Path \(2\text{-}1\text{-}3\text{-}5\), \(\ell=1\). Multiset \(=\mathit{root}_2+\mathit{root}_5-\mathit{root}_1-\varnothing=\{105,2\}+\{105,9,8\}-\{105\}=\{2,8,9,105\}\). Sorted, the 1st/2nd/3rd/4th are \(\boxed{2},\boxed{8},\boxed{9},\boxed{105}\).\\
\textbf{Query (7,8,2).} Path \(7\text{-}3\text{-}1\text{-}4\text{-}8\), \(\ell=1\). \(\mathit{root}_7=\{105,9,7\}\), \(\mathit{root}_8=\{105,3,7\}\), minus \(\mathit{root}_1=\{105\}\) gives \(\{3,7,7,9,105\}\); 2nd smallest \(\boxed{7\ }\).  Matches official output \(2,8,9,105,7\).
\end{tracebox}

\solstep{Complexity}
\begin{complexitybox}
Compression \(\bigO{N\log N}\); DFS build \(\bigO{N\log N}\); each query \(\bigO{\log N}\) LCA \(+\ \bigO{\log V}\) descent. With \(N,M\le10^5\): about \(3.4\cdot10^6\) steps -- \ac. Memory: \(\sim1.8\cdot10^6\) pool nodes of \(3\) ints \((\approx21\) MB\()\) plus the \(17\times10^5\) lifting table -- size the pool for \(\bigO{N\log N}\), never \(4N\).
\end{complexitybox}

\solstep{Verification and Edge Cases}
\begin{center}\renewcommand{\arraystretch}{1.2}
\begin{tabular}{@{}p{0.22\linewidth} p{0.30\linewidth} p{0.36\linewidth}@{}}
\toprule
\textbf{Case} & \textbf{Why tricky} & \textbf{Code behaviour}\\
\midrule
\(u=v\) & path is one node & \(\ell=u\), four roots reduce to \(\{w_u\}\); \(k=1\) \\
\(\ell=s\) (tree root) & no parent & \(\mathit{root}_0\) is the null tree; formula uniform \\
\(u\) is ancestor of \(v\) & \(\ell=u\) & sum counts spine twice, \(\mathit{root}_\ell-\mathit{root}_{p(\ell)}\) fixes \(\ell\) \\
weight \(105\), duplicates \(7,7\) & compression & ranks; equal weights share a leaf, counts add \\
\(k=j-i+1\) extreme & last descent & walks to max-rank leaf, \(\bigO{\log V}\) \\
\bottomrule
\end{tabular}\end{center}

\begin{hardeningbox}
\begin{itemize}[nosep,leftmargin=1.4em]
  \item \textbf{Four roots, not three.} Forgetting \(\mathit{root}_{\mathrm{parent}(\ell)}\) deletes \(\ell\) from the multiset; wrong whenever the answer is \(\ell\)'s value.
  \item \textbf{Carry all four through the descent.} Each level must read the left child of \(a,b,c,d\); advancing only \(a,b\) (as if two roots) silently keeps the spine.
  \item \textbf{Null root at the tree root.} \(\mathit{root}_0\) must be the shared all-zero node so that \(\ell=s\) queries do not index a missing parent.
  \item \textbf{Pool sized for \(N\log N\).} One \code{insert} per tree node; \(N=10^5\) is \(\sim1.8\cdot10^6\) nodes.
  \item \textbf{DFS depth.} A \(10^5\)-node line recurses \(10^5\) deep -- raise the stack or use an iterative DFS on judges that limit it.
\end{itemize}
\end{hardeningbox}

% ---------------------------------------------------------------
\subsection{Common Pitfalls}
% ---------------------------------------------------------------

\begin{warningbox}[Trap 1: two roots when the chains overlap]
The array case (\(\secref{c19-chairman}\)) is a difference of two \emph{nested} prefixes; a tree path is a union of two chains that meet at \(\ell\), so nested subtraction under-counts. Use the four-root rule (\(\propref{c19-fourroots}\)).
\end{warningbox}

\begin{warningbox}[Trap 2: LCA-inclusion off-by-one]
\(\ell\) belongs to the path once. \(\mathit{root}_u+\mathit{root}_v\) has it twice; subtracting \(\mathit{root}_\ell\) alone makes it zero. Only subtracting \(\mathit{root}_\ell\) \emph{and} \(\mathit{root}_{\mathrm{parent}(\ell)}\) leaves it at one.
\end{warningbox}

\begin{warningbox}[Trap 3: editing a shared root during the tree DFS]
When you \code{insert} node \(c\)'s value onto \(\mathit{root}_v\), \(\mathit{root}_v\) is shared with all other children of \(v\). Overwriting its \(\mathit{cnt}\) in place corrupts sibling paths. Clone first (\code{newnode}), write the clone -- persistence's cardinal rule (\(\secref{c19-copy}\)).
\end{warningbox}

\begin{warningbox}[Trap 4: reusing a mutable segment tree over \(tin/tout\)]
A path is \emph{not} a \(tin\)-contiguous interval (only a \emph{subtree} is). Trying to read a path off the Euler order needs HLD's chain decomposition (Ch.~22); the persistent-root route here sidesteps that entirely.
\end{warningbox}

% ---------------------------------------------------------------
\subsection{Discharging Ch.~11's deferred persistence: the persistent trie}
% ---------------------------------------------------------------

\begin{insightbox}[Same law, different alphabet -- Ch.~11's roots-by-count]
Chapter~11 Tries (\chref{c08}) built the \term{persistent trie} (\(\defref{c08-persistent}\)) to answer `how many stored numbers are \(\le x\)' or `the maximum xor with \(x\)' among a \emph{prefix} of insertions -- e.g.\ its solved LC~1707 keeps a root per prefix of `insertions,' a two-root subtraction over the \emph{bit} axis. That is the identical mechanism this section uses, with the value axis replaced by the binary trie of a key. Ch.~11 explicitly deferred the `full persistence abstraction' to this chapter (\chref{c19}); here it is: a version root that is a path copy of its predecessor, combined by difference to answer about a set of elements. Read \(\defref{c08-persistent}\)'s roots-by-count as the two-root case of \(\coderef{c19-kth}\), and \(\secref{c19-chairman}\) as the frequency-tree specialisation where the alphabet is an ordered value range. The trie is \emph{not} re-built here -- it is the same spine wearing a different alphabet.
\end{insightbox}

% ---------------------------------------------------------------
\subsection{Practice Problems}
% ---------------------------------------------------------------

\begin{stuckbox}
\begin{itemize}[nosep,leftmargin=1.4em]
  \item Did you plant \emph{one root per tree node} inside the DFS that fills the LCA lifting table (not a separate pass)?
  \item Does every path query combine \emph{four} roots \(u,v,\ell,\mathrm{parent}(\ell)\), all descending together?
  \item Is \(\mathit{root}_{\mathrm{parent}(s)}\) the null tree so that \(\ell=s\) is safe?
  \item Did you clone before writing in \code{insert}, never mutating a root shared with a sibling?
  \item Is the pool \(\bigO{N\log N}\), and is your DFS deep enough for a \(10^5\) chain?
\end{itemize}
\end{stuckbox}

\begin{practiceset}
  \practice{Arpa's letter-marked tree and Mehrdad's Dokhtar-kosh paths}{CF 741D}
    {https://codeforces.com/problemset/problem/741/D}
    {CF 2900 L5}
    {Each node carries a lowercase letter; a query path is `good' if its letters can be permuted into a palindrome (at most one letter appears an odd number of times); maximise the length of a good path. The key is a \emph{root-to-node prefix summary}: encode the root-to-\(v\) letter-parities as a bitmask (toggle a bit per letter), so the parity mask of a path \(u\text{--}v\) is \(\mathit{mask}_u\oplus\mathit{mask}_v\) corrected at \(\mathrm{lca}\) -- the same `store a per-node cumulative object, combine via the LCA' law as \(\propref{c19-fourroots}\), with a xor-mask where this section used a frequency tree. For a fixed endpoint, the best partner in a subtree is found by keeping, per mask, the deepest node seen -- small-to-large over subtrees (Ch.~28) is the common bookkeeping, but the \emph{idea} (root-to-node summary \(\to\) path answer at the LCA) is exactly \(\secref{c19-ontrees}\). Stretch \lv{5}: write the mask-prefix first, then the combine.}
  \practice{Graph and Queries}{CF 1416D}
    {https://codeforces.com/problemset/problem/1416/D}
    {CF 2600 L5}
    {A weighted graph with edge \emph{deletions}; online you are asked for the \(k\)-th largest vertex weight in the connected component of \(v\), or to delete an edge. Reverse time: deletions become insertions, and a union-find builds a \term{Kruskal reconstruction tree} whose leaves are vertices and whose internal nodes record merge time; `the component of \(v\) at time \(t\)' becomes `a subtree' of that tree. A `\(k\)-th largest weight in a subtree' is then a \(k\)-th over a \emph{static range} of the reconstruction tree's \(tin/tout\) order -- precisely \(\secref{c19-chairman}\)'s two-root descent (\(\coderef{c19-kth}\)) once you flatten, or \(\secref{c19-mst}\)'s sorted lists. The reduction chain delete\(\to\)Kruskal tree\(\to\)subtree\(\to\)Euler range\(\to\)chairman is the point: \lv{5}, but it is this chapter's machinery wearing a union-find hat.}
\end{practiceset}

\paragraph{Back-cites and forwards.}
This section is \(\secref{c19-copy}\) (path copying, \(\defref{c19-pathcopy}\)) plus \(\secref{c19-chairman}\) (the joint descent \(\coderef{c19-kth}\)), with Ch.~11's \(\mathrm{lca}\) (\chref{c11}) supplying the two extra roots and Ch.~5's \(tin/tout\) (\chref{c24}) the subtree view used only to name a range. The persistent trie Ch.~11 deferred (\(\defref{c08-persistent}\)) is discharged above. Forward: paths with \emph{updates} leave this chapter for HLD (\chref{c16}); the engine map in \(\secref{c19-sparse}\) places chairman-on-tree against the merge-sort tree and the sparse table.

% ---------------------------------------------------------------
\begin{recallbox}[Section 21.4 in 10 seconds]
Plant a persistent frequency root \(\mathit{root}_v\) on every \emph{tree node} (one \code{insert} of \(v\)'s value rank onto \(\mathit{root}_{\mathrm{parent}(v)}\), \(\bigO{N\log V}\) nodes). The path \(u\text{--}v\) is a \emph{four-root} combination \(\mathit{root}_u+\mathit{root}_v-\mathit{root}_\ell-\mathit{root}_{\mathrm{parent}(\ell)}\), \(\ell=\mathrm{lca}(u,v)\) -- \(\ell\) kept once, its ancestors deleted (\(\propref{c19-fourroots}\)). A count \(\le x\) on a path and the path's \(k\)-th smallest are each \(\bigO{\log V}\) by a four-way joint descent (\(\coderef{c19-pathkth}\)). Static only. Worked: SPOJ~COT (\(\probref{c19-cot}\), \(\coderef{c19-cot-soln}\)). Ch.~11's persistent trie (\(\defref{c08-persistent}\)) is the two-root, bit-alphabet case -- discharged, not rebuilt.
\selftest{
  \item Why does a tree path need four roots while a subarray needed two?
  \item Derive why the meeting node \(\ell\) survives exactly once under the four-root formula.
  \item On \(T\), give the path-8\text{--}10 multiset and the four-root \(k\)-th, \(k=6\), landing on \(4\).
  \item What is \(\mathit{root}_{\mathrm{parent}(s)}\) when \(\ell\) is the tree root, and why must it be the null tree?
  \item State the node/memory budget of the per-node root chain and why \(4N\) fails.
}
\end{recallbox}
